 A Interface do Usuário (UI) é responsável pela interação direta entre o operador e o sistema. Implementada como uma aplicação web responsiva, ela apresenta em tempo real os valores de temperatura e umidade coletados, disponibiliza os controles para alternância entre os modos de operação (Ligado, Desligado e Automático) e exibe alertas visuais quando os limites pré-estabelecidos são ultrapassados. A interface comunica-se com o firmware do ESP32 por meio de requisições HTTP, enviando comandos e recebendo respostas no formato JSON.

O Controle de Dados é a camada responsável por receber os dados brutos provenientes do sensor DHT11, validá-los, aplicar os procedimentos de calibração necessários e encaminhá-los às demais camadas do sistema. Essa camada implementa também a lógica de detecção de anomalias: quando um valor lido excede os limiares pré-configurados pelo usuário, o módulo aciona simultaneamente a camada de comunicação para enviar notificações por e-mail e a interface web para a exibição dos alertas e da alternância de modos.

A camada de Comunicação gerencia toda a troca de informações entre o hardware e os serviços externos. O firmware do ESP32 utiliza o protocolo HTTP/HTTPS para enviar leituras ao banco de dados remoto e para receber comandos originados na interface web.

Por fim, o Armazenamento de Dados engloba dois mecanismos complementares: o registro remoto em banco de dados relacional, acessado pela interface web para a exibição do histórico de leituras (requisito \textbf{RF06}), e o armazenamento local em arquivo \texttt{.txt} gravado na memória flash do ESP32, que garante a persistência dos dados mesmo na ausência de conectividade com a rede (requisito \textbf{RF09}).

\subsection{Descrição de Alto Nível}

\subsection{Protótipo Funcional}

\subsection{Fluxo de Dados}

O fluxo de operações do sistema pode ser descrito conforme os passos a seguir:

\begin{enumerate}

    \item O ESP32 inicializa, conecta-se à rede Wi-Fi e aguarda comandos ou eventos;
    
    \item O sensor DHT11 realiza medições periódicas de temperatura e umidade, enviando os dados ao ESP32;
    
    \item O ESP32 processa os dados recebidos e os encaminha ao banco de dados via requisição HTTP/HTTPS, registrando cada leitura;

    \item A interface web consulta o banco de dados e exibe os dados de temperatura e umidade em tempo real para o cliente, além de permitir a alternância entre os modos de operação (normal e economia de energia);

    \item Ao receber um comando de alternância de modo, o ESP32 atualiza o estado interno e aciona o LED indicador correspondente, refletindo o modo de operação ativo. 

\end{enumerate}